\documentclass[11pt]{article}
\usepackage[margin=1in]{geometry}
\usepackage{amsmath,amssymb,amsthm,hyperref}
\newtheorem{theorem}{Theorem}[section]
\newtheorem{lemma}[theorem]{Lemma}
\newtheorem{corollary}[theorem]{Corollary}
\newtheorem{remark}[theorem]{Remark}
\newcommand{\R}{\mathbb R}
\newcommand{\ip}[2]{\langle #1,#2\rangle}
\title{Optimality and Uniqueness for Net Rotation at Most One Radian}
\author{}\date{October 5, 2026 --- paper-proof draft}
\begin{document}\maketitle
\begin{abstract}
In the usual unit-width right-angled hallway, let $m(\omega)$ denote the maximum area for a prescribed net clockwise rotation $\omega$. Using fixed-angle penalized selection and first variation from the companion manuscript, we prove $m(\omega)=1+\omega^2/2$ for $0\le\omega\le1$, and identify the unique normalized optimizer. The candidate and its relaxed value are Baek's. The new steps are the fixed-angle curvature passage for every maximizing cap, an elementary arm bootstrap, and a signed-area majorization that does not require prior fan containment. A regular-closedness argument recovers arbitrary original sofas from their maximizing envelopes. This is a paper proof with explicitly identified inputs, not a Lean-checked result. Angles greater than one radian are not classified by this theorem.
\end{abstract}

\section{Statement and candidate}
Use the prescribed-net-angle convention of \texttt{paper.tex}: sofas are nonempty compact connected subsets of the plane admitting a continuous motion from the horizontal arm to the vertical arm of
\[
 L=((-\infty,1]\times[0,1])\cup([0,1]\times(-\infty,1]),
\]
with angular lift starting at zero and ending at $-\omega$. Backtracking is allowed. For $0<\omega\le1$ put
\[
 u_t=(\cos t,\sin t),\quad v_t=(-\sin t,\cos t),\quad
 c=\sec\omega-\tan\omega,\quad o=(c,1),
\]
\[
 F_\omega=\{z:z_y\ge0,\ \ip z{u_\omega}\ge0\},\qquad
 P_\omega=\{z:0\le z_y\le1,\ 0\le\ip z{u_\omega}\le1\}.
\]
Define
\[
 A(t)=o+(\omega-t)u_t-v_t,\qquad C(t)=o+t v_t-u_t,
\]
\[
 K_\omega^*=\operatorname{conv}\bigl(\{O,o\}\cup A([0,\omega])\cup C([0,\omega])\bigr),
 \qquad S_\omega^*=K_\omega^*\setminus N_\omega(K_\omega^*).
\]
The niche always uses the fan $F_\omega$, as defined in the companion manuscript. This is Baek's relaxed cap $K_{\omega,1}$ in the endpoint-strip normalization. Its supports are
\[
 p_*(t)=\omega-t+\ip o{u_t},\qquad k_*(t)=t+\ip o{v_t},
\]
and its canonical corner is
\[
 \gamma(t)=o+(\omega-t-1)u_t+(t-1)v_t.
\]

\begin{theorem}[Unrestricted fixed-angle optimum]\label{thm:main}
For $0\le\omega\le1$,
\[
 \boxed{m(\omega)=1+\frac{\omega^2}{2}.}
\]
At zero angle the only optimizer, in the fixed initial orientation and up to translation, is the unit square. For $0<\omega\le1$, every optimal sofa becomes exactly $S_\omega^*$ after the unique translation that makes its two upper endpoint supports equal to one. In particular the optimizer is unique up to congruence when no initial coordinate orientation is prescribed.
\end{theorem}

The direct feasibility, compactness and connectedness of $S_\omega^*$, as well as its exact area $1+\omega^2/2$, are proved in \texttt{relaxation.tex}. The candidate and relaxed functional are attributed there to Baek, \emph{A Conditional Upper Bound for the Moving Sofa Problem}, arXiv:2406.10725v1. The theorem here removes the restricted-class hypothesis from that note in the specified angle interval.

\section{Existence at the cap level does not need Gerver optimality}
Let $\mathcal K_\omega$ be the normalized cap class and $\mathcal A_\omega(K)=|K|-|N_\omega(K)|$.

\begin{lemma}[Compactness and upper semicontinuity]\label{lem:existence}
For $0<\omega<\pi/2$, the cap class is compact in Hausdorff distance and $\mathcal A_\omega$ is upper semicontinuous on it. If a cap of positive sofa area is available, a positive-area global cap maximizer exists.
\end{lemma}
\begin{proof}
Every cap lies in the fixed compact parallelogram $P_\omega$. Blaschke selection, or equicontinuity of its uniformly bounded support functions, gives a Hausdorff-convergent subsequence from every sequence of caps. The four pinned support values pass to the limit.

We check the defining-normal condition as well. Put
\[
 E=[0,\omega]\cup[\pi/2,\pi/2+\omega]\cup\{\pi+\omega,3\pi/2\}
\]
on the circle. The boundary-length measures of the approximating caps are supported on $E$. This follows first for finite outer approximations with normals in $E$, and then for caps by the distributional identity $\sigma_K=h_K+h_K''$. The same identity and uniform support convergence show that the limiting body's measure is supported on $E$. On each complementary normal interval $(a,b)$, of length less than $\pi$, its support solves $h''+h=0$. Hence
\[
 h(t)=\frac{\sin(b-t)}{\sin(b-a)}h(a)+\frac{\sin(t-a)}{\sin(b-a)}h(b),\quad a<t<b.
\]
The same positive-coefficient identity holds for $u_t$. Thus a point satisfying the limiting body's supporting inequalities on $E$ satisfies them on every complementary interval as well. The limiting body is therefore the intersection of its half-planes with normals in $E$, and remains a cap. This proves closedness and compactness.

If $K_j\to K$ and $z\in N_\omega(K)$, an interior angle witnesses two strict support inequalities. Uniform support convergence preserves those inequalities for all sufficiently large $j$, while the fan is fixed. Thus $\mathbf1_{N_\omega(K)}\le\liminf_j\mathbf1_{N_\omega(K_j)}$ pointwise. Fatou's lemma gives lower semicontinuity of niche area. Convex-body area is continuous under bounded Hausdorff convergence, so their difference is upper semicontinuous. It is bounded above by $|P_\omega|=\sec\omega$. A maximizing sequence therefore has a maximizing limit. A positive competitor ensures that its area functional, and consequently its convex-body area, is positive.
\end{proof}

\section{Optimality and rigidity of all maximizing caps}
\begin{lemma}[Every short-angle cap maximizer is the candidate]\label{lem:cap}
For $0<\omega\le1$,
\[
 \max_{K\in\mathcal K_\omega}\mathcal A_\omega(K)=1+\omega^2/2,
 \qquad \operatorname*{argmax}_{K\in\mathcal K_\omega}\mathcal A_\omega(K)=\{K_\omega^*\}.
\]
\end{lemma}
\begin{proof}
The feasible candidate in \texttt{relaxation.tex} has cap area $1+\omega^2/2>0$, so Lemma~\ref{lem:existence} supplies a global maximizer $K$. Apply \texttt{fixed-angle-curvature.tex} to this arbitrary maximizer, and then \texttt{arm-bootstrap.tex}. Its canonical arms satisfy
\[
 f(t)\ge1+t/2,\qquad g(t)\ge1+(\omega-t)/2.
\]
Consequently its continuously differentiable corner curve has strictly decreasing first coordinate. The theorem of \texttt{monotone-majorization.tex} now gives $\mathcal A_\omega(K)\le\mathcal A_1(K)$ without requiring prior fan containment.

The strict quadratic gap in \texttt{relaxation.tex} yields
\[
 1+\omega^2/2=\mathcal A_\omega(K_\omega^*)
 \le\mathcal A_\omega(K)\le\mathcal A_1(K)
 \le\mathcal A_1(K_\omega^*)=1+\omega^2/2.
\]
Every inequality is equality. To make the rigidity step explicit, let $f_0,g_0$ now denote the two support differences from the candidate, not the arm functions. The gap is at least
\[
 \frac12(1-\omega^2/2)\int_0^\omega f_0'^2
 +\frac12\int_0^\omega(g_0'+f_0)^2.
\]
Its first coefficient is positive. Equality forces $f_0'=0$, hence $f_0=0$ by $f_0(\omega)=0$, and then $g_0'=0$, hence $g_0=0$ by $g_0(0)=0$. The active supports coincide, so the caps coincide. Since $K$ was any maximizing cap, this classifies every cap maximizer.
\end{proof}

\section{Recovering the actual sofa, not merely its area}
\begin{lemma}[The candidate sofa is regular closed]\label{lem:regular}
For $0<\omega\le1$, $S_\omega^*=\overline{\operatorname{int}S_\omega^*}$.
\end{lemma}
\begin{proof}
The geometric lemmas of \texttt{relaxation.tex} give the convex region $D$ bounded by $O\gamma(0)$, the arc $\gamma$, and $\gamma(\omega)O$. Put $r=\omega+c-1>0$. They establish
\[
 \operatorname{int}D\subset N_\omega(K_\omega^*)\subset D,
 \quad D\subset F_\omega\cap\{\|z\|\le\sqrt2 r\},\quad \sqrt2 r<1,
\]
and $K_\omega^*$ contains the unit disk sector in $F_\omega$.

We first identify the possible points of $S_\omega^*$ on $D$. The origin belongs to the niche. For $0\le x<r$, the point $(x,0)$ belongs to a contributing quarter-plane at sufficiently small positive $t$: its first strict inequality follows from $a(0)=r>x$, where $a=p_*-1$, and its second from $-x\sin t<b(t)=k_*(t)-1>0$. Reflection gives the same assertion on the other radial edge. The niche is relatively open in $F_\omega$ and contained in $D$, so it cannot contain a point of the outer arc $\gamma$, including its endpoints: each has arbitrarily close points in the fan outside $D$. Therefore
\[
 S_\omega^*\cap D=\gamma([0,\omega]).
\]

A point of $S_\omega^*$ outside the closed set $D$ can be approximated by points of $\operatorname{int}K_\omega^*\setminus D$, since a convex body with nonempty interior is regular closed. These are interior points of the sofa. If $z=\gamma(t)$, move it a little outward along its ray from $O$. Convexity and the boundary description of $D$ show that $\lambda z\notin D$ for $\lambda>1$. For $\lambda$ sufficiently close to one, $\|\lambda z\|<1$, so $\lambda z$ is still in $K_\omega^*$. Such points again admit interior approximations outside $D$, including when the ray is a boundary ray of the fan. Letting $\lambda\downarrow1$ approximates $z$ by interior sofa points. The sofa is compact by its already proved feasibility construction, hence closed. This proves the lemma.
\end{proof}

\begin{lemma}[Closed-subset equality]
If $T$ is regular closed with finite area and $S\subset T$ is closed with $|S|=|T|$, then $S=T$.
\end{lemma}
\begin{proof}
If $z\in T\setminus S$, a ball about $z$ misses $S$. Regular closedness supplies an interior point of $T$ in that ball, and then a smaller open ball contained in $T\setminus S$. Its positive area contradicts area equality.
\end{proof}

\begin{proof}[Proof of Theorem~\ref{thm:main}]
The zero-angle case is the square theorem in \texttt{paper.tex}. Let $0<\omega\le1$. The candidate is feasible with the claimed area. For an arbitrary feasible sofa $S$, make the unique endpoint-support normalization and take its monotone envelope $T$. The fixed-angle monotonization and cap identities in the companion manuscript, Section 2, give
\[
 S\subset T=K\setminus N_\omega(K),\quad K\in\mathcal K_\omega,
 \qquad |S|\le|T|=\mathcal A_\omega(K).
\]
These identities preserve the prescribed net angle, including when the original motion backtracks. Lemma~\ref{lem:cap} bounds the last expression by $1+\omega^2/2$, proving unrestricted optimality.

If $|S|=1+\omega^2/2$, the two intervening inequalities are equalities. Thus $K$ is a global maximizing cap and Lemma~\ref{lem:cap} gives $K=K_\omega^*$. Hence $T=S_\omega^*$ and $S$ is a closed subset of $S_\omega^*$ with the same area. Lemma~\ref{lem:regular} and the closed-subset equality lemma imply $S=S_\omega^*$. No regularity assumption on the original sofa beyond closedness was added.
\end{proof}

\section{Dependency ledger and scope}
The imported companion inputs are the fixed-angle cap/monotonization identities (Section 2), penalized selection (Section 4), and floating first variation (Section 5), at commit \texttt{1ade045936f32cf76572ee668ed8aa1627772bde}. The candidate and relaxation originate in Baek's 2024 conditional-bound paper. Their concrete formulas, feasibility for $\omega\le1$, and strict equality calculation are proved in \texttt{relaxation.tex}.

The new fixed-angle curvature proof, arm bootstrap, and horizontal-monotonicity majorization are in the three correspondingly named notes. Cap attainment and recovery of arbitrary original sofas are supplied here. Neither the final Gerver optimality theorem nor its right-angle equality theorem is an input. Fan containment for an unknown maximizing corner is not assumed; in fact it follows afterwards from equality in the new majorization.

This closes the unrestricted small-angle theorem on $[0,1]$, not the full classification on $[0,\pi/2]$. The negative-results workstream tests extension beyond one radian. The proof has not been independently refereed or formalized, and no CI, Lean compilation, or TeX compilation was run.
\end{document}
